% Options for packages loaded elsewhere
\PassOptionsToPackage{unicode}{hyperref}
\PassOptionsToPackage{hyphens}{url}
\documentclass[
]{article}
\usepackage{xcolor}
\usepackage[margin=1in]{geometry}
\usepackage{amsmath,amssymb}
\setcounter{secnumdepth}{-\maxdimen} % remove section numbering
\usepackage{iftex}
\ifPDFTeX
  \usepackage[T1]{fontenc}
  \usepackage[utf8]{inputenc}
  \usepackage{textcomp} % provide euro and other symbols
\else % if luatex or xetex
  \usepackage{unicode-math} % this also loads fontspec
  \defaultfontfeatures{Scale=MatchLowercase}
  \defaultfontfeatures[\rmfamily]{Ligatures=TeX,Scale=1}
\fi
\usepackage{lmodern}
\ifPDFTeX\else
  % xetex/luatex font selection
\fi
% Use upquote if available, for straight quotes in verbatim environments
\IfFileExists{upquote.sty}{\usepackage{upquote}}{}
\IfFileExists{microtype.sty}{% use microtype if available
  \usepackage[]{microtype}
  \UseMicrotypeSet[protrusion]{basicmath} % disable protrusion for tt fonts
}{}
\makeatletter
\@ifundefined{KOMAClassName}{% if non-KOMA class
  \IfFileExists{parskip.sty}{%
    \usepackage{parskip}
  }{% else
    \setlength{\parindent}{0pt}
    \setlength{\parskip}{6pt plus 2pt minus 1pt}}
}{% if KOMA class
  \KOMAoptions{parskip=half}}
\makeatother
\usepackage{graphicx}
\makeatletter
\newsavebox\pandoc@box
\newcommand*\pandocbounded[1]{% scales image to fit in text height/width
  \sbox\pandoc@box{#1}%
  \Gscale@div\@tempa{\textheight}{\dimexpr\ht\pandoc@box+\dp\pandoc@box\relax}%
  \Gscale@div\@tempb{\linewidth}{\wd\pandoc@box}%
  \ifdim\@tempb\p@<\@tempa\p@\let\@tempa\@tempb\fi% select the smaller of both
  \ifdim\@tempa\p@<\p@\scalebox{\@tempa}{\usebox\pandoc@box}%
  \else\usebox{\pandoc@box}%
  \fi%
}
% Set default figure placement to htbp
\def\fps@figure{htbp}
\makeatother
\setlength{\emergencystretch}{3em} % prevent overfull lines
\providecommand{\tightlist}{%
  \setlength{\itemsep}{0pt}\setlength{\parskip}{0pt}}
\usepackage{graphicx}
\usepackage{float}
\usepackage{amsmath}
\usepackage{bookmark}
\IfFileExists{xurl.sty}{\usepackage{xurl}}{} % add URL line breaks if available
\urlstyle{same}
\hypersetup{
  hidelinks,
  pdfcreator={LaTeX via pandoc}}

\author{}
\date{\vspace{-2.5em}}

\begin{document}

\section{Question}\label{question}

Una nutricionista escolar, Elena, registró el consumo diario de frutas
de cada estudiante (en porciones). Los datos recopilados para 8
estudiantes se muestran en la siguiente tabla:

\includegraphics[width=3cm,height=\textheight,keepaspectratio]{tabla_datos.png}

Ricardo, un compañero de clase, analizó estos datos y afirmó: \emph{``La
mediana de este conjunto de datos es 5.75''}.

Sin embargo, Elena sospecha que Ricardo cometió un error conceptual.

\textbf{Parte 1.} ¿Cuál error conceptual cometió Ricardo al calcular la
mediana?

\#\#ANSWER1\#\#

\textbf{Parte 2.} Calcule la media aritmética correcta de los datos
(redondee a dos decimales si es necesario).

\#\#ANSWER2\#\#

\textbf{Parte 3.} Calcule la mediana correcta de los datos.

\#\#ANSWER3\#\#

\textbf{Parte 4.} Determine si la siguiente afirmación es verdadera o
falsa: \emph{``Si todos los datos de un conjunto son iguales, entonces
la media, la mediana y la moda coinciden.''}

\#\#ANSWER4\#\#

\subsection{Answerlist}\label{answerlist}

\begin{itemize}
\tightlist
\item
  Confundió la mediana con la moda
\item
  No ordenó los datos antes de buscar la mediana
\item
  Confundió la mediana con la media aritmética
\item
  Para un número par de datos, tomó solo uno de los dos valores
  centrales
\item
\item
\item
  Verdadero
\item
  Falso
\end{itemize}

\section{Solution}\label{solution}

\subsubsection{Análisis del Error (Parte
1)}\label{anuxe1lisis-del-error-parte-1}

Ricardo afirmó que la mediana es \textbf{5.75}, cuando la mediana
correcta es \textbf{5.5}.

\textbf{Error identificado:} El estudiante calculó la media aritmética
(suma de todos los datos dividida entre n) pero afirmó que ese valor era
la mediana. Este error indica que no diferencia entre las dos medidas de
tendencia central.

\textbf{Código de error:} EST-MTC-01

\textbf{Causa raíz:} No diferencia entre mediana (valor central al
ordenar) y media (promedio aritmético)

Es posible que los estudiantes que eligen otras opciones no distingan
claramente entre las tres medidas de tendencia central (media, mediana y
moda) o confundan los procedimientos de cálculo de cada una.

\subsubsection{Procedimiento Correcto (Parte 2 y
3)}\label{procedimiento-correcto-parte-2-y-3}

\textbf{Cálculo de la media aritmética:}

\textbf{Paso 1:} Sumar todos los datos.

\[\text{Suma} = 1 + 5 + 5 + 5 + 6 + 7 + 8 + 9 = 46\]

\textbf{Paso 2:} Dividir entre el número de datos.

\[\text{Media} = \frac{46}{8} = 5.75\]

\textbf{Cálculo de la mediana:}

\textbf{Paso 1:} Ordenar los datos de menor a mayor.

\[1,\; 5,\; 5,\; 5,\; 6,\; 7,\; 8,\; 9\]

\textbf{Paso 2:} Como hay 8 datos (número \textbf{par}), la mediana es
el promedio de los dos valores centrales (posiciones 4 y 5).

\[\text{Mediana} = \frac{5 + 6}{2} = 5.5\]

\textbf{Moda:} El valor 5 aparece 3 veces (la mayor frecuencia), por lo
tanto \(\text{Moda} = 5\).

\subsubsection{Propiedades de las Medidas de Tendencia Central (Parte
4)}\label{propiedades-de-las-medidas-de-tendencia-central-parte-4}

Si todos los datos de un conjunto son iguales, entonces la media, la
mediana y la moda coinciden. \(\to\) \textbf{Verdadero}

Esta afirmación es \textbf{verdadera}. La justificación se basa en las
propiedades matemáticas de las medidas de tendencia central.

\subsubsection{Caso Específico}\label{caso-especuxedfico}

En este conjunto de datos: Media = 5.75, Mediana = 5.5, Moda = 5. Las
tres medidas son diferentes, lo cual es frecuente cuando la distribución
de los datos no es simétrica.

\subsubsection{Reflexión
Metacognitiva}\label{reflexiuxf3n-metacognitiva}

Cuando identificamos el tipo de error, podemos diseñar estrategias para
evitarlo en el futuro. La autoevaluación es clave para el aprendizaje.

\subsubsection{Estrategia para Evitar el
Error}\label{estrategia-para-evitar-el-error}

\begin{enumerate}
\def\labelenumi{\arabic{enumi}.}
\tightlist
\item
  \textbf{Siempre ordenar} los datos de menor a mayor antes de calcular
  la mediana.
\item
  \textbf{Verificar} si \(n\) es par o impar para aplicar la regla
  correspondiente.
\item
  \textbf{Distinguir} claramente entre media (promedio), mediana (valor
  central) y moda (más frecuente).
\item
  \textbf{Comprobar} que la respuesta tenga sentido: la media y la
  mediana deben estar entre el valor mínimo y el máximo.
\end{enumerate}

\subsection{Answerlist}\label{answerlist-1}

\begin{itemize}
\tightlist
\item
  Incorrecto. Este no es el error que explica la respuesta del
  compañero.
\item
  Incorrecto. Este no es el error que explica la respuesta del
  compañero.
\item
  Correcto. Este es el error que cometió el compañero.
\item
  Incorrecto. Este no es el error que explica la respuesta del
  compañero.
\item
  La media correcta es \(5.75\) porciones.
\item
  La mediana correcta es \(5.5\) porciones.
\item
  Correcto. La afirmación es verdadera.
\item
  Incorrecto.
\end{itemize}

\section{Meta-information}\label{meta-information}

exname:
media\_mediana\_moda\_metacognitivo\_interpretacion\_n2\_cloze\_v1
extype: cloze exclozetype:
schoice\textbar num\textbar num\textbar schoice exsolution:
0010\textbar5.75\textbar5.5\textbar10 exshuffle: TRUE extol:
0\textbar0.01\textbar0.01\textbar0 exsection:
Estadistica-Probabilidad\textbar Pensamiento-Aleatorio\textbar Medidas-Tendencia-Central
exextra{[}Type{]}: CLOZE exextra{[}Program{]}: R exextra{[}Language{]}:
es exextra{[}Competencia{]}: Interpretacion exextra{[}Componente{]}:
Aleatorio exextra{[}Afirmacion{]}: Interpreta y diferencia las medidas
de tendencia central (media, mediana y moda) identificando errores
conceptuales en su calculo exextra{[}Evidencia{]}: Identifica errores en
el calculo de la mediana y aplica correctamente los procedimientos para
media, mediana y moda exextra{[}Nivel{]}: 2 exextra{[}Contexto{]}:
Escolar exextra{[}DOK{]}: 3 exextra{[}Bloom{]}: Analizar
exextra{[}SOLO{]}: Relacional exextra{[}TipoMetacognicion{]}:
analisis\_error

\end{document}
